% Source: intro_to_ml.tex
\begin{notechapter}
\chapter{Introduction to Machine Learning}\label{ch:intro-to-ml}
\chaptersubtitle{A summary of Chapters 1 and 2 of \textit{Machine Learning} by de Mello and Ponti, with clarifications on classifiers and $P(X,Y)$}

\section{Chapter 1: Introduction to Machine Learning}

\subsection{Key Concepts}
\textbf{Machine Learning (ML):} A field of study that enables computers to learn from data without being explicitly programmed. It focuses on developing algorithms that improve their performance with experience.

\vspace{0.5em}
\textbf{Types of Learning:}
\begin{itemize}
    \item \textbf{Supervised Learning:} The algorithm learns from labeled data, where inputs are paired with correct outputs. Examples include regression and classification.
    \item \textbf{Unsupervised Learning:} The algorithm learns patterns from unlabeled data. Examples include clustering and dimensionality reduction.
    \item \textbf{Reinforcement Learning:} The algorithm learns by interacting with an environment and receiving feedback in the form of rewards or penalties.
\end{itemize}

\vspace{0.5em}
\textbf{Statistical Learning Theory:} Provides a framework for understanding how algorithms generalize from training data to unseen data. Key concepts include bias-variance tradeoff, overfitting, and underfitting.

\vspace{0.5em}
\textbf{Model Evaluation:} Techniques like cross-validation, train-test splits, and performance metrics (e.g., accuracy, precision, recall) are used to evaluate ML models.

\subsection{Important Points}
\begin{itemize}
    \item ML algorithms aim to minimize \textbf{generalization error}, which measures how well the model performs on unseen data.
    \item The \textbf{no-free-lunch theorem} states that no single algorithm is universally better than others; the choice of algorithm depends on the problem and data.
    \item Understanding the tradeoff between \textbf{model complexity} and \textbf{generalization} is crucial to avoid overfitting or underfitting.
\end{itemize}

\section{Chapter 2: Statistical Learning Theory}

\subsection{Key Concepts}
\textbf{Learning Problem:} Formally defined as finding a function \( f: \mathcal{X} \rightarrow \mathcal{Y} \) that maps inputs \( \mathcal{X} \) to outputs \( \mathcal{Y} \) based on a dataset \( D = \{(x_1, y_1), \dots, (x_n, y_n)\} \).

\vspace{0.5em}
\textbf{Risk Minimization:} The goal is to minimize the \textbf{expected risk} \( R(f) \), which measures the expected loss of the model on unseen data:
\[
R(f) = \mathbb{E}_{(x,y) \sim P(x,y)}[L(f(x), y)],
\]
where \( L \) is a loss function (e.g., mean squared error for regression).

\vspace{0.5em}
\textbf{Empirical Risk Minimization (ERM):} Since the true distribution \( P(x,y) \) is unknown, we minimize the \textbf{empirical risk}:
\[
R_{\text{emp}}(f) = \frac{1}{n} \sum_{i=1}^n L(f(x_i), y_i).
\]

\vspace{0.5em}
\textbf{Regularization:} Introduced to prevent overfitting by adding a penalty term to the empirical risk:
\[
R_{\text{reg}}(f) = R_{\text{emp}}(f) + \lambda \Omega(f),
\]
where \( \Omega(f) \) is a regularization term (e.g., L1 or L2 norm) and \( \lambda \) controls the strength of regularization.

\subsection{Important Points}
\begin{itemize}
    \item The \textbf{Vapnik-Chervonenkis (VC) dimension} measures the capacity of a model to fit different datasets. A higher VC dimension indicates greater model complexity.
    \item The \textbf{bias-variance tradeoff} is central to understanding model performance:
    \begin{itemize}
        \item \textbf{Bias:} Error due to overly simplistic assumptions in the model.
        \item \textbf{Variance:} Error due to sensitivity to small fluctuations in the training set.
    \end{itemize}
    \item \textbf{Generalization bounds} provide theoretical guarantees on how well a model will perform on unseen data, based on its performance on the training data.
\end{itemize}

\section{Clarification: Classifiers and \( P(X, Y) \)}

A classifier is a function \( f: \mathcal{X} \rightarrow \mathcal{Y} \) that maps inputs \( \mathcal{X} \) to outputs \( \mathcal{Y} \). However, it is also related to the joint probability distribution \( P(X, Y) \). Here's how:

\subsection{Classifier as a Function}
The classifier \( f \) predicts the label \( y \in \mathcal{Y} \) for a given input \( x \in \mathcal{X} \). For example:
\begin{itemize}
    \item In binary classification, \( \mathcal{Y} = \{0, 1\} \).
    \item In multi-class classification, \( \mathcal{Y} = \{1, 2, \dots, K\} \).
\end{itemize}

\subsection{Connection to \( P(X, Y) \)}
The joint distribution \( P(X, Y) \) represents the true relationship between inputs \( X \) and outputs \( Y \). It encodes:
\begin{itemize}
    \item The probability of observing a particular input \( x \) (i.e., \( P(X = x) \)).
    \item The conditional probability of the output \( y \) given the input \( x \) (i.e., \( P(Y = y \mid X = x) \)).
\end{itemize}

The classifier \( f \) is not directly modeling \( P(X, Y) \), but it is implicitly related to it.

\subsection{Optimal Classifier and \( P(Y \mid X) \)}
The optimal classifier \( f^* \) (in terms of minimizing classification error) is the one that predicts the most likely label \( y \) for a given \( x \), based on the conditional distribution \( P(Y \mid X) \):
\[
f^*(x) = \arg\max_{y \in \mathcal{Y}} P(Y = y \mid X = x).
\]
This is known as the \textbf{Bayes classifier}.

\subsection{Learning \( f \) from Data}
In practice, we don't know \( P(Y \mid X) \) or \( P(X, Y) \). Instead, we have a dataset \( D = \{(x_1, y_1), \dots, (x_n, y_n)\} \) sampled from \( P(X, Y) \). The goal is to find a function \( f \) that approximates \( f^* \) by minimizing a loss function (e.g., cross-entropy).

\subsection{Why \( f \) Doesn't Directly Model \( P(X, Y) \)}
\begin{itemize}
    \item \( P(X, Y) \) is a joint distribution over both inputs and outputs, which is often high-dimensional and complex.
    \item The classifier \( f \) focuses only on the mapping from \( \mathcal{X} \) to \( \mathcal{Y} \), not on modeling the entire joint distribution.
    \item In many cases, we don't need to model \( P(X, Y) \) explicitly to build a good classifier. Instead, we focus on learning \( P(Y \mid X) \) or directly approximating \( f^* \).
\end{itemize}

\subsection{Example: Logistic Regression}
Logistic regression models \( P(Y \mid X) \) directly by estimating:
\[
P(Y = 1 \mid X = x) = \frac{1}{1 + e^{-(w^T x + b)}},
\]
where \( w \) and \( b \) are parameters learned from the data. The classifier \( f \) then predicts:
\[
f(x) = \begin{cases}
1 & \text{if } P(Y = 1 \mid X = x) \geq 0.5, \\
0 & \text{otherwise}.
\end{cases}
\]

\section{Bayes Classifier vs. Maximum Likelihood Classifier}

\subsection{Bayes Classifier}
The Bayes classifier is the optimal classifier in terms of minimizing the probability of misclassification. It is derived from Bayesian decision theory and makes predictions based on the posterior probability \( P(Y \mid X) \).

\begin{itemize}
    \item The Bayes classifier predicts the class \( y \) that maximizes the posterior probability:
      \[
      f^*(x) = \arg\max_{y \in \mathcal{Y}} P(Y = y \mid X = x).
      \]
    \item It is called the Bayes classifier because it is based on Bayes' theorem:
      \[
      P(Y = y \mid X = x) = \frac{P(X = x \mid Y = y) \cdot P(Y = y)}{P(X = x)}.
      \]
    \item The Bayes classifier achieves the minimum possible error rate, known as the Bayes error rate.
\end{itemize}

\subsection{Maximum Likelihood (ML) Classifier}
The Maximum Likelihood classifier is a method for estimating the parameters of a statistical model by maximizing the likelihood function. It is based on the principle of finding the parameter values that make the observed data most probable.

\begin{itemize}
    \item In classification, the ML classifier estimates the parameters of the conditional distribution \( P(X \mid Y) \) by maximizing the likelihood of the training data.
    \item Once the parameters are estimated, the ML classifier predicts the class \( y \) that maximizes the likelihood \( P(X = x \mid Y = y) \):
      \[
      \hat{y} = \arg\max_{y \in \mathcal{Y}} P(X = x \mid Y = y).
      \]
    \item The ML classifier does not directly consider the prior probabilities \( P(Y = y) \) or the posterior probabilities \( P(Y = y \mid X = x) \).
\end{itemize}

\subsection{Key Differences}
\begin{itemize}
    \item \textbf{Objective:} The Bayes classifier minimizes misclassification error, while the ML classifier maximizes the likelihood of the observed data.
    \item \textbf{Use of Priors:} The Bayes classifier incorporates prior probabilities \( P(Y = y) \), while the ML classifier does not.
    \item \textbf{Optimality:} The Bayes classifier is optimal in terms of minimizing classification error, while the ML classifier is not necessarily optimal.
\end{itemize}

\subsection{When Are They the Same?}
The Bayes classifier and the ML classifier are the same only under specific conditions:
\begin{itemize}
    \item \textbf{Equal Class Priors:} If the prior probabilities \( P(Y = y) \) are equal for all classes, the Bayes classifier reduces to the ML classifier.
    \item \textbf{Uniform Cost Function:} the Bayes classifier coincides with the ML classifier only if \emph{both} conditions hold: uniform misclassification costs (0--1 loss) \emph{and} equal class priors. With uniform costs alone, the Bayes classifier reduces to the MAP classifier $\arg\max_y P(X=x\mid Y=y)P(Y=y)$, not to the ML classifier.
\end{itemize}

\section{MAP (Maximum A Posteriori) Classifier}

\subsection{Definition}
The \textbf{MAP (Maximum A Posteriori)} classifier is another important concept in statistical learning and Bayesian decision theory. It predicts the class \( y \) that maximizes the \textbf{posterior probability} \( P(Y = y \mid X = x) \). It is derived from Bayes' theorem:

\[
P(Y = y \mid X = x) = \frac{P(X = x \mid Y = y) \cdot P(Y = y)}{P(X = x)}.
\]

Since \( P(X = x) \) is a normalizing constant that does not depend on \( y \), the MAP classifier simplifies to:

\[
f_{\text{MAP}}(x) = \arg\max_{y \in \mathcal{Y}} P(X = x \mid Y = y) \cdot P(Y = y).
\]

Here:
\begin{itemize}
    \item \( P(X = x \mid Y = y) \) is the \textbf{likelihood} of observing \( x \) given the class \( y \).
    \item \( P(Y = y) \) is the \textbf{prior probability} of class \( y \).
\end{itemize}

\subsection{Comparison with Bayes and ML Classifiers}
\begin{itemize}
    \item \textbf{Bayes Classifier:} The Bayes classifier also predicts the class \( y \) that maximizes the posterior probability \( P(Y = y \mid X = x) \). However, the Bayes classifier is optimal in terms of minimizing the expected risk (misclassification error), while the MAP classifier is optimal in terms of maximizing the posterior probability.
    \item \textbf{ML Classifier:} The ML classifier predicts the class \( y \) that maximizes the likelihood \( P(X = x \mid Y = y) \). It does not consider the prior probabilities \( P(Y = y) \), unlike the MAP classifier.
\end{itemize}

\subsection{Why Bayes and MAP Classifiers Are Often the Same in Practice}
In practice, the Bayes classifier and the MAP classifier are often the same, especially when the cost of misclassification is uniform. Here's why:
\begin{itemize}
    \item \textbf{Uniform Misclassification Cost:} When the cost of misclassifying any class is the same, the Bayes classifier reduces to maximizing the posterior probability \( P(Y = y \mid X = x) \), which is exactly what the MAP classifier does.
    \item \textbf{Prior Probabilities:} Both the Bayes classifier and the MAP classifier incorporate prior probabilities \( P(Y = y) \). When the priors are known and the misclassification costs are uniform, the two classifiers yield the same predictions.
    \item \textbf{Optimality:} The Bayes classifier is optimal in terms of minimizing the expected risk, while the MAP classifier is optimal in terms of maximizing the posterior probability. When the misclassification costs are uniform, these two objectives align, making the classifiers equivalent.
\end{itemize}

\subsection{Example}
Suppose you have a binary classification problem with:
\begin{itemize}
    \item \( P(Y = 0) = 0.7 \) (prior probability of class 0),
    \item \( P(Y = 1) = 0.3 \) (prior probability of class 1),
    \item \( P(X = x \mid Y = 0) = 0.2 \) (likelihood of \( x \) given class 0),
    \item \( P(X = x \mid Y = 1) = 0.8 \) (likelihood of \( x \) given class 1).
\end{itemize}

The MAP classifier would compute:
\[
P(Y = 0 \mid X = x) \propto 0.2 \cdot 0.7 = 0.14,
\]
\[
P(Y = 1 \mid X = x) \propto 0.8 \cdot 0.3 = 0.24.
\]
Since \( 0.24 > 0.14 \), the MAP classifier would predict \( y = 1 \).

\subsection{Key Takeaways}
\begin{itemize}
    \item The \textbf{MAP classifier} is a Bayesian approach that maximizes the posterior probability \( P(Y = y \mid X = x) \).
    \item It incorporates both the likelihood \( P(X = x \mid Y = y) \) and the prior probability \( P(Y = y) \).
    \item It is closely related to the Bayes classifier but is specifically focused on maximizing the posterior probability rather than minimizing misclassification error.
    \item The MAP classifier reduces to the ML classifier when the prior probabilities are uniform.
\end{itemize}

\section{Controlling the Cost of Misclassification}

\subsection{Cost Matrix}
The cost of misclassification is controlled using a \textbf{cost matrix} \( C \), where each element \( C_{ij} \) represents the cost of classifying an instance of class \( i \) as class \( j \). For a binary classification problem, the cost matrix might look like this:

\[
C = \begin{bmatrix}
C_{00} & C_{01} \\
C_{10} & C_{11}
\end{bmatrix},
\]

where:
\begin{itemize}
    \item \( C_{00} \): Cost of correctly classifying class \( 0 \) (usually \( 0 \)).
    \item \( C_{01} \): Cost of misclassifying class \( 0 \) as class \( 1 \).
    \item \( C_{10} \): Cost of misclassifying class \( 1 \) as class \( 0 \).
    \item \( C_{11} \): Cost of correctly classifying class \( 1 \) (usually \( 0 \)).
\end{itemize}

\subsection{Expected Risk with Cost Matrix}
The expected risk \( R(f) \) of a classifier \( f \) can be generalized to incorporate the cost matrix. Instead of simply counting misclassifications, we weight each type of error by its corresponding cost:

\[
R(f) = \mathbb{E}_{(x,y) \sim P(x,y)}[C_{y, f(x)}],
\]

where:
\begin{itemize}
    \item \( C_{y, f(x)} \) is the cost of predicting \( f(x) \) when the true class is \( y \).
    \item The expected risk is the average cost over all possible predictions, weighted by the joint distribution \( P(x, y) \).
\end{itemize}

\subsection{Optimal Classifier with Cost Matrix}
The optimal classifier \( f^* \) minimizes the expected risk \( R(f) \). For a given input \( x \), the classifier predicts the class \( y \) that minimizes the expected cost:

\[
f^*(x) = \arg\min_{y \in \mathcal{Y}} \sum_{y' \in \mathcal{Y}} P(Y = y' \mid X = x) \cdot C_{y', y}.
\]

This is a generalization of the Bayes classifier that incorporates the cost matrix.

\subsection{Uniform Cost of Misclassification}
When the cost of misclassification is \textbf{uniform}, all misclassification errors are equally costly. For example, in a binary classification problem, the cost matrix might look like this:

\[
C = \begin{bmatrix}
0 & 1 \\
1 & 0
\end{bmatrix}.
\]

Here:
\begin{itemize}
    \item \( C_{01} = 1 \): Cost of misclassifying class \( 0 \) as class \( 1 \).
    \item \( C_{10} = 1 \): Cost of misclassifying class \( 1 \) as class \( 0 \).
\end{itemize}

In this case, the optimal classifier reduces to the standard Bayes classifier, which maximizes the posterior probability \( P(Y = y \mid X = x) \).

\subsection{Non-Uniform Cost of Misclassification}
When the cost of misclassification is \textbf{non-uniform}, the optimal classifier takes into account the different costs. For example, if misclassifying class \( 1 \) as class \( 0 \) is more costly than the reverse, the classifier will predict class \( 1 \) more readily (it becomes more conservative in predicting class \( 0 \)).

\subsubsection{Example: Medical Diagnosis}
In a medical diagnosis scenario:
\begin{itemize}
    \item \( C_{01} = 1 \): Cost of misclassifying a healthy patient as sick.
    \item \( C_{10} = 10 \): Cost of misclassifying a sick patient as healthy.
\end{itemize}

The cost matrix would be:
\[
C = \begin{bmatrix}
0 & 1 \\
10 & 0
\end{bmatrix}.
\]

\subsubsection{Optimal Classifier Behavior}
Using the medical diagnosis cost matrix:
\[
C = \begin{bmatrix}
0 & 1 \\
10 & 0
\end{bmatrix},
\]
the optimal classifier predicts class \( 1 \) (sick) when the expected cost of predicting \( 0 \), \( 10\,P(Y = 1 \mid X = x) \), exceeds the expected cost of predicting \( 1 \), \( 1\cdot P(Y = 0 \mid X = x) \), i.e.\ when \( P(Y = 1 \mid X = x) > 1/11 \approx 0.09 \) instead of \( 0.5 \). The classifier is thus \emph{less} conservative in predicting class \( 1 \): it reduces false negatives at the expense of more false positives.

\subsection{Implementing Cost-Sensitive Classification}
In practice, cost-sensitive classification can be implemented in several ways:
\begin{itemize}
    \item \textbf{Weighted Loss Functions:} Modify the loss function to include the cost matrix. For example, in logistic regression, the cross-entropy loss can be weighted by the costs.
    \item \textbf{Resampling:} Adjust the training data to reflect the costs. For example, oversample the class with higher misclassification costs.
    \item \textbf{Threshold Adjustment:} Adjust the decision threshold of a classifier to account for the costs. For example, in binary classification, the threshold for predicting class \( 1 \) can be lowered if the cost of false negatives is high.
    \item \textbf{Cost-Sensitive Algorithms:} Use algorithms specifically designed to handle cost matrices, such as cost-sensitive decision trees or cost-sensitive support vector machines.
\end{itemize}

\subsection{Example: Threshold Adjustment}
Suppose we have a binary classifier that outputs probabilities \( P(Y = 1 \mid X = x) \). Normally, we predict class \( 1 \) if \( P(Y = 1 \mid X = x) \geq 0.5 \). However, if the cost of false negatives is high, we might lower the threshold to \( 0.3 \):

\[
f(x) = \begin{cases}
1 & \text{if } P(Y = 1 \mid X = x) \geq 0.3, \\
0 & \text{otherwise}.
\end{cases}
\]

This reduces the likelihood of misclassifying class \( 1 \) as class \( 0 \), at the expense of increasing false positives.

\subsection{Key Takeaways}
\begin{itemize}
    \item The cost of misclassification is controlled using a \textbf{cost matrix}, which assigns specific penalties to different types of errors.
    \item The \textbf{expected risk} of a classifier is generalized to incorporate the cost matrix, and the optimal classifier minimizes this risk.
    \item When the cost of misclassification is \textbf{uniform}, the optimal classifier reduces to the standard Bayes classifier.
    \item When the cost is \textbf{non-uniform}, the classifier must account for the different costs, often by adjusting the decision threshold or using cost-sensitive algorithms.
\end{itemize}

\section{Die-Throwing Example from the Book}

\subsection{Setup}
In the die-throwing example from the book:
\begin{itemize}
    \item \( X \): A binary random variable where \( X = 1 \) if an even number occurs and \( X = 0 \) otherwise.
    \item \( Y \): A binary random variable where \( Y = 1 \) if the number is prime and \( Y = 0 \) otherwise.
    \item The joint probability \( P(X, Y) \) is given based on the outcomes of throwing a die.
    \item The algorithm converges to a classifier \( g \) with the following induced joint probabilities:
      \[
      P_g(X = 0, Y = 0) = 0, \quad P_g(X = 0, Y = 1) = \frac{3}{6},
      \]
      \[
      P_g(X = 1, Y = 0) = \frac{3}{6}, \quad P_g(X = 1, Y = 1) = 0.
      \]
\end{itemize}

\subsection{Expected Risk \( R(g) \)}
The expected risk \( R(g) \) measures how well the classifier \( g \) approximates the true relationship between \( X \) and \( Y \). It is computed as the sum of the absolute differences between the true joint probabilities \( P(X, Y) \) and the induced joint probabilities \( P_g(X, Y) \):
\[
R(g) = \sum_{x \in \{0,1\}} \sum_{y \in \{0,1\}} \left| P(X = x, Y = y) - P_g(X = x, Y = y) \right|.
\]

\subsection{Step-by-Step Computation}
Assume the true joint probabilities \( P(X, Y) \) are as follows:
\[
P(X = 0, Y = 0) = \frac{1}{6}, \quad P(X = 0, Y = 1) = \frac{2}{6},
\]
\[
P(X = 1, Y = 0) = \frac{2}{6}, \quad P(X = 1, Y = 1) = \frac{1}{6}.
\]

The expected risk \( R(g) \) is computed as:
\[
R(g) = \left| \frac{1}{6} - 0 \right| + \left| \frac{2}{6} - \frac{3}{6} \right| + \left| \frac{2}{6} - \frac{3}{6} \right| + \left| \frac{1}{6} - 0 \right|.
\]
Simplifying:
\[
R(g) = \frac{1}{6} + \frac{1}{6} + \frac{1}{6} + \frac{1}{6} = \frac{4}{6} = \frac{2}{3}.
\]

\subsection{Connection to the General Formula for \( R(g) \)}
The computation of \( R(g) \) in this example is a specific instance of the general expected risk formula:
\[
R(g) = \mathbb{E}_{(X,Y) \sim P(X,Y)}[L(g(X), Y)],
\]
where one may be tempted to define the ``loss'' as:
\[
L(g) = \left| P(X = x, Y = y) - P_g(X = x, Y = y) \right|.
\]
This measures the discrepancy between the true distribution and the distribution induced by \( g \).

\begin{remark}[This is not an expected loss]
The quantity above is \emph{not} of the form \( \mathbb{E}[L(g(X), Y)] \): a loss function compares a
prediction \( g(x) \) with a label \( y \), whereas \( |P(X = x, Y = y) - P_g(X = x, Y = y)| \) compares two
distributions. It is the \( \ell^1 \) distance between \( P(X,Y) \) and \( P_g(X,Y) \) (twice their total
variation distance). For a deterministic classifier, however, it is closely related to the 0--1 risk:
\( P_g(X = x, Y = g(x)) = P(X = x) \) and \( P_g(X = x, Y = y) = 0 \) for \( y \neq g(x) \), so
\[
\sum_{x,y} \left| P(x, y) - P_g(x, y) \right|
= \sum_x \Big( \big(P(x) - P(x, g(x))\big) + \sum_{y \neq g(x)} P(x, y) \Big)
= 2\, P\big(g(X) \neq Y\big).
\]
Here \( P(g(X)\neq Y) = \frac{1}{3} \), consistent with \( \frac{2}{3} \) above.
\end{remark}

\subsection{Bayes Classifier \( f^* \)}
Since we know the true joint distribution \( P(X, Y) \), we can derive the
\textbf{Bayes classifier \( f^* \)}. The Bayes classifier predicts the most likely label \( y \) for a given \( x \), based on the conditional distribution \( P(Y \mid X) \):
\[
f^*(x) = \arg\max_{y \in \{0,1\}} P(Y = y \mid X = x).
\]

For this example:
\begin{itemize}
    \item If \( X = 0 \):
      \[
      P(Y = 0 \mid X = 0) = \frac{P(X = 0, Y = 0)}{P(X = 0)} = \frac{\frac{1}{6}}{\frac{1}{6} + \frac{2}{6}} = \frac{1}{3},
      \]
      \[
      P(Y = 1 \mid X = 0) = \frac{P(X = 0, Y = 1)}{P(X = 0)} = \frac{\frac{2}{6}}{\frac{1}{6} + \frac{2}{6}} = \frac{2}{3}.
      \]
      Thus, \( f^*(0) = 1 \) (since \( P(Y = 1 \mid X = 0) > P(Y = 0 \mid X = 0) \)).

    \item If \( X = 1 \):
      \[
      P(Y = 0 \mid X = 1) = \frac{P(X = 1, Y = 0)}{P(X = 1)} = \frac{\frac{2}{6}}{\frac{2}{6} + \frac{1}{6}} = \frac{2}{3},
      \]
      \[
      P(Y = 1 \mid X = 1) = \frac{P(X = 1, Y = 1)}{P(X = 1)} = \frac{\frac{1}{6}}{\frac{2}{6} + \frac{1}{6}} = \frac{1}{3}.
      \]
      Thus, \( f^*(1) = 0 \) (since \( P(Y = 0 \mid X = 1) > P(Y = 1 \mid X = 1) \)).
\end{itemize}

\subsection{Expected Risk of the Bayes Classifier \( R(f^*) \)}
The Bayes classifier \( f^* \) is optimal, meaning it achieves the minimum possible expected risk. However, this does not imply that \( R(f^*) = 0 \). Instead, \( R(f^*) \) represents the \textbf{Bayes error rate}, which is the minimum achievable error rate given the uncertainty in \( Y \) conditioned on \( X \).

\subsubsection{Why \( R(f^*) \neq 0 \) in General}
The Bayes classifier \( f^* \) predicts the most likely label \( y \) for a given \( x \) based on the conditional distribution \( P(Y \mid X) \):
\[
f^*(x) = \arg\max_{y \in \{0,1\}} P(Y = y \mid X = x).
\]
However, if \( P(Y \mid X) \) is not deterministic (i.e., there is uncertainty in \( Y \) given \( X \)), the Bayes classifier will still make errors. Specifically, for a given \( x \), the probability of error is:
\[
P(\text{error} \mid X = x) = 1 - \max_{y \in \{0,1\}} P(Y = y \mid X = x).
\]
The expected risk \( R(f^*) \) is then the average of these error probabilities over all \( x \), weighted by \( P(X = x) \):
\[
R(f^*) = \mathbb{E}_{X \sim P(X)}[P(\text{error} \mid X)].
\]

\subsubsection{Connection to the General Expected Risk}
The two definitions of expected risk are connected as follows:

\begin{itemize}
    \item The general definition of expected risk \( R(f) \) applies to any classifier \( f \) and any loss function \( L \):
      \[
      R(f) = \mathbb{E}_{(x,y) \sim P(x,y)}[L(f(x), y)].
      \]
      Here, \( L(f(x), y) \) quantifies the cost of predicting \( f(x) \) when the true label is \( y \).

    \item For the Bayes classifier \( f^* \), the loss function \( L(f(x), y) \) is often defined as the \textbf{0-1 loss} (also called the misclassification error):
      \[
      L(f(x), y) = \mathbb{I}(f(x) \neq y),
      \]
      where \( \mathbb{I}(\cdot) \) is the indicator function that returns \( 1 \) if the prediction \( f(x) \) is incorrect and \( 0 \) otherwise. With this loss function, the expected risk becomes:
      \[
      R(f^*) = \mathbb{E}_{(x,y) \sim P(x,y)}[\mathbb{I}(f^*(x) \neq y)].
      \]
      This represents the probability of misclassification, averaged over all possible inputs \( x \) and outputs \( y \).

    \item By the \textbf{law of total expectation}, the expected risk can be rewritten as:
      \[
      R(f^*) = \mathbb{E}_{X \sim P(X)}[\mathbb{E}_{Y \mid X}[\mathbb{I}(f^*(x) \neq y) \mid X]].
      \]
      Here, the inner expectation \( \mathbb{E}_{Y \mid X}[\mathbb{I}(f^*(x) \neq y) \mid X] \) is the \textbf{conditional probability of error} given \( X = x \). To see why, note that:
      \[
      \mathbb{E}_{Y \mid X}[\mathbb{I}(f^*(x) \neq y) \mid X = x] = \sum_{y \in \mathcal{Y}} \mathbb{I}(f^*(x) \neq y) \cdot P(Y = y \mid X = x).
      \]
      The indicator function \( \mathbb{I}(f^*(x) \neq y) \) is \( 1 \) when \( f^*(x) \neq y \) and \( 0 \) otherwise. Therefore, the expectation simplifies to:
      \[
      \mathbb{E}_{Y \mid X}[\mathbb{I}(f^*(x) \neq y) \mid X = x] = \sum_{y \neq f^*(x)} P(Y = y \mid X = x).
      \]
      This sum is exactly the probability that \( Y \neq f^*(x) \) given \( X = x \), which is the conditional probability of error:
      \[
      P(\text{error} \mid X = x) = P(Y \neq f^*(x) \mid X = x).
      \]

    \item Therefore, the expected risk of the Bayes classifier \( f^* \) can be expressed as:
      \[
      R(f^*) = \mathbb{E}_{X \sim P(X)}[P(\text{error} \mid X)].
      \]
      This shows that \( R(f^*) \) is the average of the conditional error probabilities \( P(\text{error} \mid X) \), weighted by the probability \( P(X = x) \).
\end{itemize}

\subsubsection{Key Takeaway}
The two definitions of expected risk are equivalent in the context of classification when:
\begin{itemize}
    \item The loss function \( L(f(x), y) \) is the 0-1 loss \( \mathbb{I}(f(x) \neq y) \).
    \item The classifier \( f \) is the Bayes classifier \( f^* \), which minimizes the conditional error probability \( P(\text{error} \mid X) \).
\end{itemize}
Thus, the expression:
\[
R(f^*) = \mathbb{E}_{X \sim P(X)}[P(\text{error} \mid X)],
\]
is a specific instance of the general expected risk formula, tailored to the Bayes classifier and the 0-1 loss.

\subsubsection{Step-by-Step Calculation for the Die-Throwing Example}
For the die-throwing example, let's compute \( R(f^*) \) explicitly.

\begin{itemize}
    \item For \( X = 0 \):
      \[
      P(Y = 0 \mid X = 0) = \frac{1}{3}, \quad P(Y = 1 \mid X = 0) = \frac{2}{3}.
      \]
      The Bayes classifier predicts \( f^*(0) = 1 \), and the probability of error is:
      \[
      P(\text{error} \mid X = 0) = 1 - \max\left( \frac{1}{3}, \frac{2}{3} \right) = 1 - \frac{2}{3} = \frac{1}{3}.
      \]

    \item For \( X = 1 \):
      \[
      P(Y = 0 \mid X = 1) = \frac{2}{3}, \quad P(Y = 1 \mid X = 1) = \frac{1}{3}.
      \]
      The Bayes classifier predicts \( f^*(1) = 0 \), and the probability of error is:
      \[
      P(\text{error} \mid X = 1) = 1 - \max\left( \frac{2}{3}, \frac{1}{3} \right) = 1 - \frac{2}{3} = \frac{1}{3}.
      \]

    \item The overall expected risk \( R(f^*) \) is the weighted average of the error probabilities:
      \[
      R(f^*) = P(X = 0) \cdot P(\text{error} \mid X = 0) + P(X = 1) \cdot P(\text{error} \mid X = 1).
      \]
      Substituting the values:
      \[
      R(f^*) = \left( \frac{1}{6} + \frac{2}{6} \right) \cdot \frac{1}{3} + \left( \frac{2}{6} + \frac{1}{6} \right) \cdot \frac{1}{3} = \frac{3}{6} \cdot \frac{1}{3} + \frac{3}{6} \cdot \frac{1}{3} = \frac{1}{6} + \frac{1}{6} = \frac{1}{3}.
      \]
\end{itemize}

\subsubsection{Interpretation}
\begin{itemize}
    \item The Bayes classifier \( f^* \) achieves the minimum possible expected risk, but this risk is not necessarily zero. In this example, \( R(f^*) = \frac{1}{3} \), which reflects the inherent uncertainty in \( Y \) given \( X \).
    \item With the 0--1 loss, the classifier \( g \) of the book (\( g(0) = 1 \), \( g(1) = 0 \)) is in fact \emph{equal} to the Bayes classifier \( f^* \), and its misclassification risk is \( P(g(X) \neq Y) = P(X=0,Y=0) + P(X=1,Y=1) = \frac{1}{3} = R(f^*) \). The value \( \frac{2}{3} \) obtained above is not a misclassification risk: as shown in the remark following its computation, it is twice the misclassification risk, so the two numbers cannot be compared.
\end{itemize}

\subsubsection{Key Takeaways}
\begin{itemize}
    \item The Bayes classifier \( f^* \) minimizes the expected risk, but \( R(f^*) \) is not necessarily zero. It represents the \textbf{Bayes error rate}, which quantifies the inherent uncertainty in \( Y \) given \( X \).
    \item In this example, \( R(f^*) = \frac{1}{3} \), and the classifier \( g \) found by the algorithm is optimal: \( P(g(X)\neq Y) = \frac{1}{3} \). (The quantity \( \frac{2}{3} \) is \( 2P(g(X)\neq Y) \).)
    \item The computation of \( R(f^*) \) involves calculating the probability of error for each \( x \) and averaging these probabilities weighted by \( P(X = x) \).
\end{itemize}

\subsubsection{Relationship Between \( R(g) \) and \( R(f^*) \)}
For any classifier \( g \), the expected risk \( R(g) \) satisfies:
\[
R(g) \geq R(f^*),
\]
where \( R(f^*) \) is the Bayes error rate. However, the relationship between \( R(g) \) and \( R(f^*) \) depends on the specific classifier \( g \):

\begin{itemize}
    \item \textbf{Strict Inequality \( R(g) > R(f^*) \):}
      This occurs when \( g \) is suboptimal compared to \( f^* \). For example, if \( g \) makes predictions that are less accurate than \( f^* \) for some inputs \( x \), then \( R(g) \) will be larger than \( R(f^*) \).

    \item \textbf{Equality \( R(g) = R(f^*) \):}
      This occurs in the following cases:
      \begin{itemize}
          \item \( g \) is equivalent to \( f^* \), i.e., \( g(x) = f^*(x) \) for all \( x \) where \( P(X = x) > 0 \).
          \item \( g \) differs from \( f^* \) only on inputs \( x \) where \( P(X = x) = 0 \).
          \item \( g \) is another Bayes-optimal classifier, e.g., when two classes are equally likely given \( X = x \).
      \end{itemize}
\end{itemize}

\subsubsection{Key Takeaway}
While \( R(g) \geq R(f^*) \) always holds, the strict inequality \( R(g) > R(f^*) \) is not guaranteed for all classifiers \( g \neq f^* \). Equality \( R(g) = R(f^*) \) is possible if \( g \) is equivalent to \( f^* \) or differs from \( f^* \) only on inputs with zero probability.

\subsubsection{Dependence on the Loss Function}
The relationship \( R(g) \geq R(f^*) \) holds for any loss function, but the definition of the Bayes classifier \( f^* \) depends on the choice of loss function. For example:

\begin{itemize}
    \item \textbf{0-1 Loss:}
      For the 0-1 loss:
      \[
      L(f(x), y) = \mathbb{I}(f(x) \neq y),
      \]
      the Bayes classifier \( f^* \) minimizes the probability of misclassification:
      \[
      f^*(x) = \arg\max_{y \in \mathcal{Y}} P(Y = y \mid X = x).
      \]
      In this case, \( R(f^*) \) is the Bayes error rate, and any other classifier \( g \) satisfies:
      \[
      R(g) \geq R(f^*).
      \]

    \item \textbf{Cross-Entropy Loss:}
      For the cross-entropy loss:
      \[
      L(f(x), y) = -\log f_y(x),
      \]
      the Bayes classifier \( f^* \) minimizes the expected cross-entropy risk:
      \[
      R(f) = \mathbb{E}_{(x,y) \sim P(x,y)}[-\log f_y(x)],
      \]
      Here \( f(x) = (f_1(x), \dots, f_K(x)) \) is a predicted probability distribution over the classes, and the optimal \( f^* \) predicts the true conditional distribution, \( f^*_y(x) = P(Y = y \mid X = x) \) (by Gibbs' inequality), not just the most likely class. For any other classifier \( g \), we still have:
      \[
      R(g) \geq R(f^*).
      \]

    \item \textbf{Squared Error Loss:}
      For the squared error loss (commonly used in regression):
      \[
      L(f(x), y) = (f(x) - y)^2,
      \]
      the Bayes classifier \( f^* \) minimizes the expected squared error risk:
      \[
      R(f) = \mathbb{E}_{(x,y) \sim P(x,y)}[(f(x) - y)^2].
      \]
      In this case, the Bayes classifier \( f^* \) is the \textbf{conditional expectation}:
      \[
      f^*(x) = \mathbb{E}[Y \mid X = x].
      \]
      This is because the conditional expectation minimizes the expected
      squared error
      (see \hyperref[intro-to-ml:sec:squared_error_loss]{detailed explanations below}). For any other regressor \( g \), we have:
      \[
      R(g) \geq R(f^*).
      \]
\end{itemize}

\subsubsection{Key Takeaway}
The relationship \( R(g) \geq R(f^*) \) holds for any loss function, but the definition of the Bayes classifier \( f^* \) depends on the choice of loss function:
\begin{itemize}
    \item For the 0-1 loss, \( f^* \) minimizes the probability of misclassification.
    \item For cross-entropy loss, \( f^* \) minimizes the expected cross-entropy risk.
    \item For squared error loss, \( f^* \) is the conditional expectation \( \mathbb{E}[Y \mid X = x] \).
\end{itemize}
In all cases, the Bayes classifier \( f^* \) achieves the minimum possible expected risk \( R(f^*) \), and any other classifier \( g \) satisfies \( R(g) \geq R(f^*) \).

\subsubsection{Details about the Squared Error Loss}\label{intro-to-ml:sec:squared_error_loss}

For the squared error loss (commonly used in regression):
\[
L(f(x), y) = (f(x) - y)^2,
\]
the Bayes classifier \( f^* \) minimizes the expected squared error risk:
\[
R(f) = \mathbb{E}_{(x,y) \sim P(x,y)}[(f(x) - y)^2].
\]

\subsubsection{Decomposing the Expected Squared Error}
To minimize \( R(f) \), we can use the law of total expectation to decompose the global expected squared error into a weighted average of conditional expectations:
\[
R(f) = \mathbb{E}_{(x,y) \sim P(x,y)}[(f(x) - y)^2] = \mathbb{E}_{X}[\mathbb{E}_{Y \mid X}[(f(X) - Y)^2 \mid X]].
\]
Here:
\begin{itemize}
    \item The outer expectation \( \mathbb{E}_{X} \) is taken over the marginal distribution \( P(X) \).
    \item The inner expectation \( \mathbb{E}_{Y \mid X} \) is taken over the conditional distribution \( P(Y \mid X) \).
\end{itemize}

\subsubsection{Conditional Expected Squared Error}
For a fixed \( x \), the conditional expected squared error is:
\[
\mathbb{E}[(f(x) - Y)^2 \mid X = x].
\]
This represents the average squared error for a fixed input \( x \), where the average is taken over the conditional distribution \( P(Y \mid X = x) \).

\subsubsection{Minimizing the Conditional Expected Squared Error}
To minimize the global expected squared error \( R(f) \), we must minimize the conditional expected squared error \( \mathbb{E}[(f(x) - Y)^2 \mid X = x] \) for each \( x \). This is because:
\[
R(f) = \sum_{x \in \mathcal{X}} P(X = x) \cdot \mathbb{E}[(f(x) - Y)^2 \mid X = x],
\]
where:
\begin{itemize}
    \item \( P(X = x) \) is the probability of \( X = x \), and
    \item \( \mathbb{E}[(f(x) - Y)^2 \mid X = x] \) is the conditional expected squared error for a fixed \( x \).
\end{itemize}
Since \( P(X = x) \) are positive weights and \( \mathbb{E}[(f(x) - Y)^2 \mid X = x] \) are non-negative terms, the only way to minimize \( R(f) \) is to minimize each term \( \mathbb{E}[(f(x) - Y)^2 \mid X = x] \) individually. If any term were not minimized, it would contribute a positive amount to \( R(f) \), preventing \( R(f) \) from being minimized.

\subsubsection{Why the Bayes Classifier is the Conditional Expectation}
The Bayes classifier \( f^* \) for squared error loss is the conditional expectation:
\[
f^*(x) = \mathbb{E}[Y \mid X = x].
\]
This is because the conditional expectation minimizes the conditional expected squared error. Let's break this down step by step.

\subsubsection{Step 1: Expanding the Squared Term}
For a fixed \( x \), the conditional expected squared error is:
\[
\mathbb{E}[(f(x) - Y)^2 \mid X = x] = \mathbb{E}[f(x)^2 - 2f(x)Y + Y^2 \mid X = x].
\]
Using the linearity of expectation, this becomes:
\[
f(x)^2 - 2f(x)\mathbb{E}[Y \mid X = x] + \mathbb{E}[Y^2 \mid X = x].
\]

\subsubsection{Step 2: Minimizing with Respect to \( f(x) \)}
To find the value of \( f(x) \) that minimizes this expression, we take the derivative with respect to \( f(x) \) and set it to zero:
\[
\frac{d}{df(x)} \left( f(x)^2 - 2f(x)\mathbb{E}[Y \mid X = x] + \mathbb{E}[Y^2 \mid X = x] \right) = 2f(x) - 2\mathbb{E}[Y \mid X = x] = 0.
\]
Solving for \( f(x) \), we get:
\[
f(x) = \mathbb{E}[Y \mid X = x].
\]

\subsubsection{Step 3: Uniqueness of the Minimizer}

To prove that \( f^*(x) = \mathbb{E}[Y \mid X = x] \) is the unique minimizer of the conditional expected squared error, consider any other function \( g(x) \neq f^*(x) \). For a fixed \( x \), the conditional expected squared error for \( g(x) \) is:

\[
\mathbb{E}[(g(x) - Y)^2 \mid X = x].
\]

Let's expand this expression:

\[
\mathbb{E}[(g(x) - Y)^2 \mid X = x] = \mathbb{E}[g(x)^2 - 2g(x)Y + Y^2 \mid X = x].
\]

Using the linearity of expectation, this becomes:

\[
g(x)^2 - 2g(x)\mathbb{E}[Y \mid X = x] + \mathbb{E}[Y^2 \mid X = x].
\]

Recall that \( f^*(x) = \mathbb{E}[Y \mid X = x] \), so we can substitute \( f^*(x) \) for \( \mathbb{E}[Y \mid X = x] \):

\[
g(x)^2 - 2g(x)f^*(x) + \mathbb{E}[Y^2 \mid X = x].
\]

Now, let's compare this to the conditional expected squared error for \( f^*(x) \):

\[
\mathbb{E}[(f^*(x) - Y)^2 \mid X = x] = f^*(x)^2 - 2f^*(x)\mathbb{E}[Y \mid X = x] + \mathbb{E}[Y^2 \mid X = x].
\]

Again, substituting \( f^*(x) = \mathbb{E}[Y \mid X = x] \), this simplifies to:

\[
f^*(x)^2 - 2f^*(x)f^*(x) + \mathbb{E}[Y^2 \mid X = x] = f^*(x)^2 - 2f^*(x)^2 + \mathbb{E}[Y^2 \mid X = x] = -f^*(x)^2 + \mathbb{E}[Y^2 \mid X = x].
\]

Now, subtract the conditional expected squared error for \( f^*(x) \) from that of \( g(x) \):

\begin{align*}
&\mathbb{E}[(g(x) - Y)^2 \mid X = x] - \mathbb{E}[(f^*(x) - Y)^2 \mid X = x] \\
&\qquad= \left( g(x)^2 - 2g(x)f^*(x) + \mathbb{E}[Y^2 \mid X = x] \right) - \left( -f^*(x)^2 + \mathbb{E}[Y^2 \mid X = x] \right).
\end{align*}

Simplify the right-hand side:

\[
= g(x)^2 - 2g(x)f^*(x) + \mathbb{E}[Y^2 \mid X = x] + f^*(x)^2 - \mathbb{E}[Y^2 \mid X = x].
\]

Notice that the \( \mathbb{E}[Y^2 \mid X = x] \) terms cancel out:

\[
= g(x)^2 - 2g(x)f^*(x) + f^*(x)^2.
\]

This expression is a perfect square:

\[
= (g(x) - f^*(x))^2.
\]

Thus, we have shown that:

\[
\mathbb{E}[(g(x) - Y)^2 \mid X = x] - \mathbb{E}[(f^*(x) - Y)^2 \mid X = x] = (g(x) - f^*(x))^2.
\]

\subsubsection{Interpretation of the Result}

This equation tells us that the difference in the conditional expected squared errors between \( g(x) \) and \( f^*(x) \) is exactly \( (g(x) - f^*(x))^2 \). Since \( (g(x) - f^*(x))^2 \) is always non-negative and equals zero only when \( g(x) = f^*(x) \), it follows that:

\begin{itemize}
    \item If \( g(x) \neq f^*(x) \), then \( \mathbb{E}[(g(x) - Y)^2 \mid X = x] > \mathbb{E}[(f^*(x) - Y)^2 \mid X = x] \).
    \item If \( g(x) = f^*(x) \), then \( \mathbb{E}[(g(x) - Y)^2 \mid X = x] = \mathbb{E}[(f^*(x) - Y)^2 \mid X = x] \).
\end{itemize}

\subsubsection{Conclusion}

This proves that \( f^*(x) = \mathbb{E}[Y \mid X = x] \) is the unique minimizer of the conditional expected squared error. Any other function \( g(x) \neq f^*(x) \) will result in a larger expected squared error, and the amount of this increase is exactly \( (g(x) - f^*(x))^2 \). Therefore,
\( f^*(x) \) is the optimal predictor under squared error loss.

\subsubsection{Step 4: Global Uniqueness}
Since \( f^*(x) = \mathbb{E}[Y \mid X = x] \) is the unique minimizer of the conditional expected squared error for every \( x \), it is also the unique minimizer of the global expected squared error \( R(f) \).

\subsubsection{Interpretation}
The conditional expectation \( \mathbb{E}[Y \mid X = x] \) is the "center of mass" of the distribution of \( Y \) given \( X = x \). It is the value that minimizes the expected squared error because:
\begin{itemize}
    \item It balances the deviations of \( Y \) from \( f(x) \) in a way that minimizes the average squared distance.
    \item Any other value of \( f(x) \) would result in a larger expected squared error.
\end{itemize}

\subsubsection{Bayes Optimality}
The Bayes classifier \( f^* \) is defined as the function that minimizes the expected risk for a given loss function. For squared error loss, the Bayes classifier is:
\[
f^*(x) = \mathbb{E}[Y \mid X = x].
\]
This is because \( f^*(x) \) achieves the minimum possible expected squared error risk:
\[
R(f^*) = \mathbb{E}_{(x,y) \sim P(x,y)}[(f^*(x) - y)^2].
\]

\subsubsection{Key Takeaway}
The Bayes classifier for squared error loss is the conditional expectation \( \mathbb{E}[Y \mid X = x] \). This is because the conditional expectation is the unique minimizer of the expected squared error, making it the optimal predictor under squared error loss. For any other regressor \( g \), we have:
\[
R(g) \geq R(f^*).
\]

\subsection{Empirical Risk vs. Average Error}

In the context of machine learning, \textbf{Empirical Risk} and \textbf{average error} are closely related concepts, but they are not exactly the same thing. Let's clarify the distinction between the two, especially in the context of the book \textit{Machine Learning} by de Mello and Ponti.

\subsubsection{Empirical Risk}

Empirical Risk is a concept from statistical learning theory. It refers to the average loss (or error) of a model over the \textbf{training dataset}. Formally, for a given loss function \( L \) and a dataset \( D = \{(x_1, y_1), \dots, (x_n, y_n)\} \), the empirical risk estimate \( R_{\text{emp}}(f) \) of a model \( f \) is defined as:

\[
R_{\text{emp}}(f) = \frac{1}{n} \sum_{i=1}^n L(f(x_i), y_i),
\]

where:
\begin{itemize}
    \item \( L(f(x_i), y_i) \) is the loss (or error) for the \( i \)-th training example,
    \item \( n \) is the number of training examples.
\end{itemize}

The empirical risk measures how well the model \( f \) performs on the \textbf{training data}. It is called "empirical" because it is computed based on the observed data (the training set), rather than the true underlying distribution of the data.

\subsubsection{Average Error}

The term \textbf{average error} is more general and can refer to the average loss (or error) of a model over \textbf{any dataset}, not just the training set. For example:
\begin{itemize}
    \item If the dataset is the \textbf{training set}, the average error is the same as the empirical risk.
    \item If the dataset is the \textbf{test set} or a \textbf{validation set}, the average error refers to the model's performance on unseen data.
\end{itemize}

Formally, for a dataset \( D' = \{(x'_1, y'_1), \dots, (x'_m, y'_m)\} \), the average error \( R_{\text{avg}}(f) \) of a model \( f \) is:

\[
R_{\text{avg}}(f) = \frac{1}{m} \sum_{i=1}^m L(f(x'_i), y'_i),
\]

where:
\begin{itemize}
    \item \( L(f(x'_i), y'_i) \) is the loss for the \( i \)-th example in the dataset \( D' \),
    \item \( m \) is the number of examples in \( D' \).
\end{itemize}

\subsubsection{Key Differences}

\begin{itemize}
    \item \textbf{Scope of Application}:
        \begin{itemize}
            \item \textbf{Empirical Risk}: Specifically refers to the average loss over the \textbf{training dataset}.
            \item \textbf{Average Error}: Can refer to the average loss over \textbf{any dataset}, including the training set, test set, or validation set.
        \end{itemize}

    \item \textbf{Purpose}:
        \begin{itemize}
            \item \textbf{Empirical Risk}: Used in the context of \textbf{training a model}. The goal of many machine learning algorithms is to minimize the empirical risk (e.g., via Empirical Risk Minimization, or ERM).
            \item \textbf{Average Error}: Used to evaluate the model's performance on \textbf{any dataset}, including unseen data. It is often used to assess generalization performance.
        \end{itemize}

    \item \textbf{Theoretical vs. Practical}:
        \begin{itemize}
            \item \textbf{Empirical Risk}: Has a strong theoretical foundation in statistical learning theory. It is used to approximate the \textbf{true risk} (expected loss over the entire data distribution) when the true distribution is unknown.
            \item \textbf{Average Error}: A more practical term used to describe the model's performance on a specific dataset.
        \end{itemize}
\end{itemize}

\subsubsection{Connection to the Book}

In \textit{Machine Learning} by de Mello and Ponti, the authors likely discuss \textbf{Empirical Risk} in the context of \textbf{Empirical Risk Minimization (ERM)}, which is a fundamental principle in machine learning. ERM aims to find the model \( f \) that minimizes the empirical risk:

\[
f^* = \arg\min_f R_{\text{emp}}(f).
\]

The authors may also discuss \textbf{average error} in the context of evaluating a model's performance on a test set or validation set. For example, the average error on a test set is often used to estimate the model's generalization performance.

\subsubsection{When Are They the Same?}

\begin{itemize}
    \item If the dataset used to compute the average error is the \textbf{training set}, then the average error is the same as the empirical risk.
    \item However, if the dataset is a \textbf{test set} or \textbf{validation set}, the average error is not the same as the empirical risk, as it measures performance on unseen data rather than the training data.
\end{itemize}

\subsubsection{Summary}

\begin{itemize}
    \item \textbf{Empirical Risk}: Average loss over the \textbf{training set}. Used during model training.
    \item \textbf{Average Error}: Average loss over \textbf{any dataset} (training, test, or validation). Used to evaluate model performance.
\end{itemize}

In the context of de Mello and Ponti's book, the two terms are related but not identical. Empirical Risk is a specific case of average error, applied to the training set, while average error is a more general term that can apply to any dataset.

\subsection{Empirical Risk as an Estimator}

While the empirical risk is often computed as an average error over the training set, it is important to understand its role as an estimator of the true risk in statistical learning theory. This leads us to the concept of the empirical risk as an estimator.

\subsubsection{True Risk (Expected Risk)}

The true risk \( R(f) \) is the expected loss of a model \( f \) over the entire data distribution \( P(X, Y) \). It is defined as:

\[
R(f) = \mathbb{E}_{(X,Y) \sim P(X,Y)}[L(f(X), Y)],
\]

where:
\begin{itemize}
    \item \( L(f(X), Y) \) is the loss function,
    \item \( \mathbb{E} \) denotes the expectation over the true data distribution \( P(X, Y) \).
\end{itemize}

The true risk is a \textbf{fixed quantity} for a given model \( f \) and data distribution \( P(X, Y) \).

\subsubsection{Empirical Risk as an Estimator}

Since we typically do not have access to the true data distribution \( P(X, Y) \), we use the \textbf{training dataset} \( D = \{(x_1, y_1), \dots, (x_n, y_n)\} \) to approximate the true risk. The empirical risk estimator \( R_{\text{emp}}(f) \) is defined as:

\[
R_{\text{emp}}(f) = \frac{1}{n} \sum_{i=1}^n L(f(X_i), Y_i),
\]

where \( \{(X_1, Y_1), \dots, (X_n, Y_n)\} \) is a \textbf{random sample} from \( P(X, Y) \). Here:
\begin{itemize}
    \item \( (X_i, Y_i) \) are \textbf{random variables}, representing the process of sampling from \( P(X, Y) \).
    \item \( R_{\text{emp}}(f) \) is a \textbf{random variable} because it depends on the random sample \( \{(X_1, Y_1), \dots, (X_n, Y_n)\} \).
\end{itemize}

The empirical risk \( R_{\text{emp}}(f) \) is an unbiased estimator of the true risk \( R(f) \) because:
\[
\mathbb{E}[R_{\text{emp}}(f)] = R(f),
\]
where the expectation is taken over all possible realizations of the random sample \( \{(X_1, Y_1), \dots, (X_n, Y_n)\} \).

\subsubsection{Empirical Risk as an Estimate}

For a specific realization of the random sample (i.e., the observed dataset \( D = \{(x_1, y_1), \dots, (x_n, y_n)\} \)), the empirical risk is computed as:

\[
R_{\text{emp}}(f) = \frac{1}{n} \sum_{i=1}^n L(f(x_i), y_i),
\]

where \( (x_i, y_i) \) are the \textbf{observed values} of \( (X_i, Y_i) \). This is a \textbf{fixed value} for the given dataset \( D \).

\subsubsection{Key Distinction: Estimator vs. Estimate}

\begin{itemize}
    \item \textbf{Estimator}: The formula \( R_{\text{emp}}(f) = \frac{1}{n} \sum_{i=1}^n L(f(X_i), Y_i) \) is an \textbf{estimator}. It is a function of random variables and is itself a random variable.
    \item \textbf{Estimate}: The value \( R_{\text{emp}}(f) = \frac{1}{n} \sum_{i=1}^n L(f(x_i), y_i) \) is an \textbf{estimate}. It is a fixed value computed from the observed dataset \( D \).
\end{itemize}

\subsubsection{Why Does This Matter?}

This distinction is crucial for understanding the \textbf{statistical properties} of machine learning models:
\begin{itemize}
    \item When analyzing the behavior of the empirical risk (e.g., its expectation, variance, or performance as an estimator of the true risk), we use the \textbf{estimator} \( R_{\text{emp}}(f) = \frac{1}{n} \sum_{i=1}^n L(f(X_i), Y_i) \).
    \item When computing the empirical risk for a specific dataset, we use the \textbf{estimate} \( R_{\text{emp}}(f) = \frac{1}{n} \sum_{i=1}^n L(f(x_i), y_i) \).
\end{itemize}

\subsection{What is a Random Sample?}

A \textbf{random sample} refers to a collection of \textbf{random variables} drawn from a probability distribution. In the context of machine learning and statistics:

\subsubsection{Random Sample as Random Variables}

A random sample is a set of random variables \( \{(X_1, Y_1), \dots, (X_n, Y_n)\} \), where each \( (X_i, Y_i) \) is an independent and identically distributed (i.i.d.) draw from the true data distribution \( P(X, Y) \). Here:
\begin{itemize}
    \item \( (X_i, Y_i) \) are \textbf{random variables}, meaning they are not fixed values but rather represent potential outcomes governed by \( P(X, Y) \).
    \item The notation \( (X_i, Y_i) \sim P(X, Y) \) means that each \( (X_i, Y_i) \) is drawn from the distribution \( P(X, Y) \).
\end{itemize}

This set of random variables \( \{(X_1, Y_1), \dots, (X_n, Y_n)\} \) is called a \textbf{random sample} because it represents the process of sampling from \( P(X, Y) \).

\subsubsection{Observed Dataset}

Once the random sample \( \{(X_1, Y_1), \dots, (X_n, Y_n)\} \) is realized (i.e., the random variables take on specific values), we obtain the \textbf{observed dataset} \( D = \{(x_1, y_1), \dots, (x_n, y_n)\} \). Here:
\begin{itemize}
    \item \( (x_i, y_i) \) are \textbf{deterministic values} (observed data points) that result from the random variables \( (X_i, Y_i) \) taking on specific realizations.
    \item The dataset \( D \) is fixed once observed, but it is the result of a random process (the sampling from \( P(X, Y) \)).
\end{itemize}

\subsubsection{Key Distinction}

\begin{itemize}
    \item \textbf{Random Sample}: Refers to the set of \textbf{random variables} \( \{(X_1, Y_1), \dots, (X_n, Y_n)\} \) before they are observed. These are not fixed values but rather represent the process of sampling from \( P(X, Y) \).
    \item \textbf{Observed Dataset}: Refers to the set of \textbf{deterministic values} \( \{(x_1, y_1), \dots, (x_n, y_n)\} \) obtained after the random variables \( (X_i, Y_i) \) are realized. This dataset is fixed once observed.
\end{itemize}

\subsubsection{Why Does This Matter?}

This distinction is crucial for understanding the statistical foundations of machine learning:
\begin{itemize}
    \item \textbf{Randomness in Sampling}:
        \begin{itemize}
            \item The randomness comes from the process of sampling \( \{(X_1, Y_1), \dots, (X_n, Y_n)\} \) from \( P(X, Y) \).
            \item The observed dataset \( D = \{(x_1, y_1), \dots, (x_n, y_n)\} \) is just one possible realization of this random process.
        \end{itemize}

    \item \textbf{Empirical Risk as a Random Variable}:
        \begin{itemize}
            \item The empirical risk \( R_{\text{emp}}(f) \) is a function of the dataset \( D \), which is a realization of the random sample \( \{(X_1, Y_1), \dots, (X_n, Y_n)\} \).
            \item Because \( D \) is a random sample, \( R_{\text{emp}}(f) \) is a \textbf{random variable} whose value depends on the specific realization of \( D \).
        \end{itemize}

    \item \textbf{Statistical Inference}:
        \begin{itemize}
            \item When we analyze the properties of \( R_{\text{emp}}(f) \) (e.g., its expectation or variance), we are analyzing its behavior across different possible realizations of the random sample \( \{(X_1, Y_1), \dots, (X_n, Y_n)\} \).
        \end{itemize}
\end{itemize}
\end{notechapter}
